\chapterbackmatter{Solutions to Weinberg Exercises}
\ExerciseEditorialNote

\begin{enumerate}
\WeinbergSolution{1}
Choose the gauge in which the static external field is time independent and
write the charged Klein--Gordon field as
\[
 \phi(\mathbf x,t)=\sum_N\left[
 a_Nu_N(\mathbf x)e^{-\ii E_Nt}
 {}+b_N^\dagger v_N(\mathbf x)e^{\ii E_Nt}\right].
\]
(A sum over \(N\) includes continuum integration.)  For a purely
electrostatic field define
\[
 \epsilon_N^{(+)}(\mathbf x)=E_N+e\mathcal A^0(\mathbf x),\qquad
 \epsilon_N^{(-)}(\mathbf x)=E_N-e\mathcal A^0(\mathbf x).
\]
The equal-time commutators \([\phi(\mathbf x),\phi^\dagger(\mathbf y)]=0\)
and \([\phi(\mathbf x),\pi(\mathbf y)]=\ii\delta^3(\mathbf x-\mathbf y)\)
give the two completeness identities
\begin{align*}
 \sum_N\{u_N(\mathbf x)u_N^*(\mathbf y)
            -v_N(\mathbf x)v_N^*(\mathbf y)\}&=0,\\
 \sum_N\{u_N(\mathbf x)\epsilon_N^{(+)}(\mathbf y)u_N^*(\mathbf y)
            +v_N(\mathbf x)\epsilon_N^{(-)}(\mathbf y)v_N^*(\mathbf y)\}
 &=\delta^3(\mathbf x-\mathbf y).
\end{align*}
Consequently every square-integrable test function has the expansion
\[
 \boxed{f(\mathbf x)=\sum_N[c_Nu_N(\mathbf x)+d_Nv_N(\mathbf x)]}
\]
with
\[
 \boxed{\begin{aligned}
 c_N&=\int \dd^3y\,u_N^*(\mathbf y)
       \epsilon_N^{(+)}(\mathbf y)f(\mathbf y),\\
 d_N&=\int \dd^3y\,v_N^*(\mathbf y)
       \epsilon_N^{(-)}(\mathbf y)f(\mathbf y).
 \end{aligned}}
\]
These are the charged Klein--Gordon inner-product projections.  If a static
vector potential is also present, the same formulas hold because only the
covariant time derivative enters the equal-time product.

\WeinbergSolution{2}
Time-translation invariance permits the energy transform
\[
 \Pi_{\mathcal A}^*(\mathbf x,\mathbf y;E)
 =\int \dd\tau\,e^{\ii E\tau}
   \Pi_{\mathcal A}^*(x^0,\mathbf x;x^0-\tau,\mathbf y).
\]
Near an unperturbed one-particle pole, insert the mode expansion from
Problem~1 on both sides of the Dyson equation.  A single proper self-energy
insertion produces a double pole.  Comparing its coefficient with
\[
 \frac1{E-E_N-\delta E_N}
 =\frac1{E-E_N}+\frac{\delta E_N}{(E-E_N)^2}+\cdots
\]
gives
\[
 \boxed{\displaystyle
 \delta E_N=-\int \dd^3x\,\dd^3y\,
 u_N^*(\mathbf x)\,
 \Pi_{\mathcal A}^*(\mathbf x,\mathbf y;E_N)\,
 u_N(\mathbf y) .}
\]
The Klein--Gordon normalization used in Problem~1 already contains the pole
residue, so no additional \(1/(2E_N)\) occurs.  In a degenerate subspace the
right-hand side is a matrix that must be diagonalized.  Its real part shifts
the level, while an unstable level has
\(\Gamma_N=-2\,\operatorname{Im}\delta E_N\).

\WeinbergSolution{3}
The imaginary part of Eq.~(14.3.45) gives the electric-dipole rate
\[
 \Gamma_N=\frac{4\alpha}{3}\sum_{E_M<E_N}
 (E_N-E_M)|\mathbf v_{MN}|^2
 =\frac{4\alpha}{3}\sum_{E_M<E_N}
 (E_N-E_M)^3|\mathbf x_{MN}|^2 .
\]
For a \(2p\) state the only lower E1 level is \(1s\), and
\[
 \omega_{21}=\frac{3}{8}\alpha^2m_e,\qquad
 \int_0^\infty \dd r\,r^3R_{10}(r)R_{21}(r)
 =\frac{256}{81\sqrt6}\,a_0 .
\]
The angular integral gives, for each magnetic sublevel,
\[
 \sum_i|\langle1s|x^i|2p,m\rangle|^2
 =\frac13\left(\frac{256a_0}{81\sqrt6}\right)^2
 =\frac{32768}{59049}a_0^2 .
\]
It follows that
\[
 \boxed{\displaystyle
 \Gamma_{2p\to1s}=\frac{2^8}{3^8}\alpha^5m_e .}
\]
Numerically this is \(6.27\times10^8\ {\rm s}^{-1}\), or
\(\tau_{2p}\simeq1.60\ {\rm ns}\).  Fine-structure and recoil corrections
are beyond the requested order.

\WeinbergSolution{4}
Let \(H\) be the Schrödinger--Coulomb Hamiltonian and
\(\Delta_{MN}=E_M-E_N\).  After the free on-shell mass and residue
subtractions, old-fashioned perturbation theory gives the state-dependent
part of the scalar self-energy as
\[
 -\frac{g^2}{2}\int\frac{\dd^3k}{(2\pi)^3\omega_k}
 \left\langle N\left|
 e^{\ii\mathbf k\cdot\mathbf x}
 \frac1{H-E_N+\omega_k}
 e^{-\ii\mathbf k\cdot\mathbf x}
 \right|N\right\rangle_{\rm sub},
 \qquad \omega_k=\sqrt{\mathbf k^2+m_\phi^2}.
\]
In the stated range \(\Delta_{MN}\ll m_\phi\ll Z\alpha m_e\), expand the
resolvent in \(H-E_N\) and use isotropy of the \(1s\) state.  The constant
and linear terms are precisely the on-shell mass and residue subtractions.
The first surviving term is
\[
 \delta E_N=-\frac{g^2\langle\mathbf p^2\rangle_N}{6m_e^2}
 \int\frac{\dd^3k}{(2\pi)^3}
 \frac{\mathbf k^2}{(\mathbf k^2+m_\phi^2)^2}.
\]
Dimensional regularization isolates the nonanalytic, scheme-independent
term,
\[
 \int\frac{\dd^3k}{(2\pi)^3}
 \frac{\mathbf k^2}{(\mathbf k^2+m_\phi^2)^2}
 =-\frac{3m_\phi}{8\pi}.
\]
Since \(\langle\mathbf p^2\rangle_{1s}=(Z\alpha m_e)^2\),
\[
 \boxed{\displaystyle
 \delta E_{1s}=\frac{g^2}{16\pi}(Z\alpha)^2m_\phi
 {}+O\!\left(g^2(Z\alpha)^4m_e,\,
 g^2(Z\alpha)^2\frac{m_\phi^2}{Z\alpha m_e}\right).}
\]
The positive sign means that the level is made less deeply bound.  Terms
analytic in \(m_\phi^2\) belong to the shorter-distance matching indicated
in the remainder.

\WeinbergSolution{5}
The expansion used in Problem~4 is not uniform when \(m_\phi=0\).  After
the on-shell mass and residue pieces have been removed, the dipole term
\(e^{\ii\mathbf k\cdot\mathbf x}=1+\ii\mathbf k\cdot\mathbf x+\cdots\)
leaves
\[
 \delta E_N^{\rm soft}
 =\frac{g^2}{12\pi^2}\sum_M
 \Delta_{MN}^3|\mathbf x_{MN}|^2
 \ln\frac{\Lambda}{|\Delta_{MN}|}
 +\hbox{nonlogarithmic matching terms}.
\]
The artificial scale \(\Lambda\) cancels against the hard part.  For the
ground state define the scalar Bethe mean excitation energy by
\[
 \ln\Delta E_{1s}^{(\phi)}
 =\frac{\sum_M\Delta_{M1}|\mathbf v_{M1}|^2
                  \ln\Delta_{M1}}
        {\sum_M\Delta_{M1}|\mathbf v_{M1}|^2}.
\]
The double-commutator sum rule gives
\[
 \sum_M\Delta_{M1}|\mathbf v_{M1}|^2
 =\frac{Ze^2}{2m_e^2}|\psi_{1s}(0)|^2
 =2(Z\alpha)^4m_e .
\]
Thus, to the leading logarithmic accuracy fixed without a full hard
one-loop matching calculation,
\[
 \delta E_{1s}(m_\phi=0)
 =\frac{g^2}{6\pi^2}(Z\alpha)^4m_e
 \left[\ln\frac{m_e}{\Delta E_{1s}^{(\phi)}}+O(1)\right].
\]
Because \(\Delta E_{1s}^{(\phi)}\) is of atomic order, the leading
dependence may be displayed more simply as
\[
 \boxed{\displaystyle
 \delta E_{1s}(0)=\frac{g^2}{6\pi^2}(Z\alpha)^4m_e
 \left[\ln\frac1{(Z\alpha)^2}+O(1)\right].}
\]
This also explains why simply setting \(m_\phi=0\) in the leading answer to
Problem~4 is invalid: the linear term disappears and atomic excitation
energies themselves regulate the infrared region.
\end{enumerate}
